\section*{Chapitre \ref{ch:g4:changes-of-state} --- Fondre et bouillir~: les changements d'état}

\begin{solution}{exo:g4:changes-of-state:1}
La fusion~; la solidification (congélation)~; l'ébullition~; la
condensation.
\end{solution}

\begin{solution}{exo:g4:changes-of-state:2}
\omterm{def:g4:changes-of-state:melting-point}{Température de fusion} \qty{0}{\celsius}, \omterm{def:g4:changes-of-state:boiling-point}{température d'ébullition}
\qty{100}{\celsius}. Le chocolat fond vers \qty{34}{\celsius} --- la
\omterm{def:g3:temperature-thermometers:temperature}{température} de la bouche.
\end{solution}

\begin{solution}{exo:g4:changes-of-state:3}
\qty{0}{\celsius} --- et toujours \qty{0}{\celsius} cinq minutes plus
tard~: tant que la glace et l'eau cohabitent, la \omterm{def:g3:temperature-thermometers:temperature}{température} est
clouée à la \omterm{def:g4:changes-of-state:melting-point}{température de fusion}.
\end{solution}

\begin{solution}{exo:g4:changes-of-state:4}
Pendant qu'une matière fond ou bout, sa \omterm{def:g3:temperature-thermometers:temperature}{température} reste immobile au
repère. La chaleur est entièrement dépensée pour le changement d'état
lui-même --- transformer le \omterm{def:g3:solids-liquids-gases:liquid}{liquide} en vapeur --- et non pour monter.
\end{solution}

\begin{solution}{exo:g4:changes-of-state:5}
\omterm{def:g3:solids-liquids-gases:solid}{Solide} à \qty{20}{\celsius} (au-dessous de $60$)~; \omterm{def:g3:solids-liquids-gases:liquid}{liquide} à
\qty{80}{\celsius} (au-dessus de $60$)~; encore \omterm{def:g3:solids-liquids-gases:solid}{solide} à
\qty{35}{\celsius} --- alors qu'une tablette de chocolat dans une
poche y aurait déjà fondu.
\end{solution}

\begin{solution}{exo:g4:changes-of-state:6}
Pas plus chaude~: les deux casseroles sont à \qty{100}{\celsius}, le
palier d'ébullition. La chaleur supplémentaire ne fait que fabriquer
de la vapeur plus vite --- du \omterm{def:g3:solids-liquids-gases:gas}{gaz} gaspillé, les mêmes pâtes.
\end{solution}

\begin{solution}{exo:g4:changes-of-state:7}
Aux deux tronçons plats --- les paliers. Partout où la courbe reste
horizontale alors que la cuisinière continue de chauffer, un
changement d'état est en cours.
\end{solution}

\begin{solution}{exo:g4:changes-of-state:8}
L'or ($1064 < 1200$) \omterm{def:g1:floating-and-sinking:words}{coule}~; le fer ($1538 > 1200$) reste \omterm{def:g3:solids-liquids-gases:solid}{solide}. Le
four manque le repère du fer de $1538 - 1200 = 338$ degrés.
\end{solution}

\begin{solution}{exo:g4:changes-of-state:9}
La glace qui fond cloue la boisson à \qty{0}{\celsius}~: toute la
chaleur qui arrive sert à faire \omterm{def:g2:ice-water-steam:meltfreeze}{fondre} la glace, et rien à réchauffer
la boisson. Ce n'est qu'une fois le dernier éclat disparu que la
\omterm{def:g3:temperature-thermometers:temperature}{température} peut se mettre à monter.
\end{solution}

\begin{solution}{exo:g4:changes-of-state:10}
Sa meilleure hypothèse~: la matière est de la glace --- de l'eau
gelée. La loi qui fait des repères des marques d'identité~: les
\omterm{def:g4:changes-of-state:melting-point}{températures de fusion} et d'ébullition d'une matière pure ne changent
jamais~; retrouver exactement les deux nombres de l'eau vaut donc une
signature.
\end{solution}

\begin{solution}{exo:g4:changes-of-state:11}
La certitude que l'eau bout toujours à \qty{100}{\celsius}. Très haut
sur un sommet, elle bout bien au-dessous --- chaude, mais pas à
\qty{100}{\celsius} --- si bien que l'eau «~bouillante~» est trop
tiède pour faire infuser le thé. Le coupable est l'\omterm{def:g2:air-around-us:air}{air} raréfié, qui
appuie moins sur l'eau, comme le montrera le chapitre sur la pression
de l'\omterm{def:g2:air-around-us:air}{air}.
\end{solution}
